% ============================================================
\section{\ours{}: Methodology}
\label{sec:method}
% ============================================================

Figure~\ref{fig:architecture} illustrates the end-to-end \ours{} pipeline.
It consists of two stages: (1) protocol-constrained hard-negative generation
and (2) supervised contrastive fine-tuning of Moirai.

% --- Figure 1: Architecture ---
\begin{figure}[t]
  \centering
  % TikZ System Architecture Diagram for Figure 1
% Compile within paper/main.tex (requires tikz, arrows.meta, positioning, shapes)

\begin{tikzpicture}[
  node distance=0.4cm and 0.7cm,
  box/.style={draw, rounded corners=2pt, minimum width=1.5cm, minimum height=0.6cm,
              align=center, font=\footnotesize, inner sep=4pt},
  stage/.style={draw, dashed, rounded corners=4pt, inner sep=5pt, fill=gray!5},
  arrow/.style={-{Stealth[length=2.5pt]}, semithick},
  label/.style={font=\scriptsize\itshape},
]

%% ---- Stage 1: Stealth-Controlled Negative Generation ----
\node[box, fill=green!15] (ciciot)    {CICIoT2023};
\node[box, fill=green!15, right=of ciciot] (benign)   {Benign $\mathbf{X}_b$};
\node[box, fill=orange!20, right=0.9cm of benign] (diffts) {Archetype\\Perturbation\\$P_k(\mathbf{X}_b,s)$};
\node[box, fill=orange!20, below=0.4cm of diffts] (constraint) {Protocol\\Constraint\\$\mathcal{C}_\pi$};
\node[box, fill=orange!20, right=0.8cm of diffts] (hardneg) {Stealth-Ctrl\\Neg. $\mathbf{X}_a$};

\draw[arrow] (ciciot) -- (benign);
\draw[arrow] (benign) -- (diffts);
\draw[arrow] (diffts) -- (hardneg);
\draw[arrow] (diffts) -- (constraint);
\draw[arrow] (constraint.east) to[bend right=30] node[label,right]{retry if fail} (diffts.south);

%% ---- Data Mixture ----
\node[box, fill=yellow!15, right=0.8cm of hardneg] (mix) {50/50\\Data Mix};

\draw[arrow] (hardneg) -- (mix) node[midway,above,label]{50\%};
\draw[arrow] (benign.north) to[bend left=35] (mix) node[pos=0.7,above,label]{50\%};

%% ---- Stage 2: Fine-Tuning ----
\node[box, fill=blue!15, right=0.8cm of mix] (moirai) {Moirai\\(fine-tuned)};
\node[box, fill=blue!10, above=0.35cm of moirai] (projhead) {Projection\\Head $h_\phi$};
\node[box, fill=blue!10, below=0.35cm of moirai] (nll)    {$\mathcal{L}_\text{NLL}$};
\node[box, fill=blue!10, right=0.8cm of projhead] (supcon) {$\mathcal{L}_\text{SupCon}$};
\node[box, fill=blue!20, right=0.8cm of nll]     (combined) {$\mathcal{L}=\mathcal{L}_\text{NLL}$\\$+\lambda\mathcal{L}_\text{SupCon}$};

\draw[arrow] (mix) -- (moirai);
\draw[arrow] (moirai) -- (projhead);
\draw[arrow] (moirai) -- (nll);
\draw[arrow] (projhead) -- (supcon);
\draw[arrow] (supcon) -- (combined);
\draw[arrow] (nll) -- (combined);

%% ---- Inference ----
\node[box, fill=purple!15, right=0.8cm of combined] (detect) {Anomaly\\Score $\hat{s}$};
\node[box, fill=purple!20, right=0.7cm of detect]  (output) {Attack / Benign};

\draw[arrow] (combined) -- (detect) node[midway,above,label]{CI scoring};
\draw[arrow] (detect) -- (output) node[midway,above,label]{$\hat{s} > \delta$};

%% ---- Stage bounding boxes ----
\begin{scope}[on background layer]
  \node[stage, fit=(diffts)(constraint)(hardneg), label=above:{\textbf{Stage 1: Generation}}] {};
  \node[stage, fit=(moirai)(projhead)(nll)(supcon)(combined), label=above:{\textbf{Stage 2: Fine-Tuning}}] {};
\end{scope}

\end{tikzpicture}

  \caption{
    \ours{} architecture.
    \textbf{Stage 1}: Diffusion-TS generates hard-negative attack sequences
    conditioned on benign traffic and an attack pattern $P_k$, with a
    constraint-retry loop ensuring protocol validity.
    \textbf{Stage 2}: Moirai is fine-tuned on a 70/20/10 mixture of
    benign, hard-negative, and real-attack sequences using a combined
    NLL + supervised contrastive loss.
    At inference, out-of-confidence-interval scoring detects anomalies.
  }
  \label{fig:architecture}
\end{figure}

% ============================================================
\subsection{Stage 1: Protocol-Constrained Hard-Negative Generation}
\label{sec:method:gen}
% ============================================================

\textbf{Diffusion-TS backbone.}
We adopt Diffusion-TS~\cite{zhang2024diffusionts}, which models the
score function of a multivariate time series as a sum of trend and
seasonality components.  Given a benign sample $\mathbf{X}_b$, we
apply a forward diffusion process to obtain a noisy latent, then
guide the reverse process toward the target attack pattern $P_k$
at stealth level $s$:
\begin{equation}
  \mathbf{X}_a = \text{DiffusionTS}\!\left(\mathbf{X}_b,\, P_k,\, s\right).
\end{equation}
The trend component uses a Savitzky-Golay filter, and the seasonality
component retains the top-5 FFT frequencies, enabling interpretable
attack injection that preserves global traffic statistics.

\textbf{Protocol constraint validation.}
Raw generated samples may violate IoT protocol specifications (e.g.,
packet sizes outside the Modbus 7--260 byte range).  Algorithm~\ref{alg:hng}
describes the constraint-retry loop used during generation.

\begin{algorithm}[h]
\caption{Protocol-Constrained Hard-Negative Generation}
\label{alg:hng}
\begin{algorithmic}[1]
\Require Benign sequence $\mathbf{X}_b$, attack pattern $P_k$,
         stealth level $s$, protocol $\pi$, max retries $R=3$
\Ensure Hard-negative $\mathbf{X}_a$ satisfying $\mathcal{C}_\pi$
\For{$r = 1, \ldots, R$}
  \State $\mathbf{X}_a \leftarrow \text{DiffusionTS}(\mathbf{X}_b, P_k, s)$
  \State $\text{report} \leftarrow \mathcal{C}_\pi(\mathbf{X}_a)$
  \If{$\text{report.is\_valid}$}
    \State \Return $\mathbf{X}_a$
  \EndIf
  \State $s \leftarrow \min(s + 0.01,\; 1.0)$  \Comment{Relax stealth slightly}
\EndFor
\State \Return $\mathbf{X}_a$  \Comment{Return best attempt with warning}
\end{algorithmic}
\end{algorithm}

We implement three IoT protocol validators:
\textbf{Modbus} (packet size 7--260 bytes, timing 1--1000 ms, valid function codes),
\textbf{MQTT} (packet size 2--1024 bytes, QoS levels 0--2, keep-alive 60--300 s),
and \textbf{CoAP} (packet size 4--1280 bytes, valid message types, token length 0--8).
Each validator operates at three strictness tiers (strict / moderate / permissive)
to support different deployment requirements.

% ============================================================
\subsection{Stage 2: Moirai Fine-Tuning with Supervised Contrastive Loss}
\label{sec:method:finetune}
% ============================================================

\textbf{Moirai for anomaly scoring.}
Moirai~\cite{woo2024moirai} is a universal Transformer forecaster
that produces probabilistic predictions via mixture distributions.
Given context window $\mathbf{X}_{1:96}$, it predicts a distribution
over $\mathbf{X}_{97:128}$.  We use the $1-\alpha$ confidence interval
(CI) at $\alpha = 0.05$ as an anomaly signal: timestep $t$ is flagged
if the observed value falls outside the predicted CI.  The sample-level
anomaly score is the fraction of flagged timesteps:
\begin{equation}
  \hat{s}(\mathbf{X}) = \frac{1}{T'}\sum_{t=97}^{128}
    \mathbf{1}\!\left[x_t \notin \text{CI}_{0.95}(\hat{X}_t)\right].
\end{equation}
A sample is classified as an attack if $\hat{s}(\mathbf{X}) > \delta$,
where $\delta$ is calibrated on a validation set
(Section~\ref{sec:experiments:setup}).

\textbf{Gradient fix for fine-tuning.}
The \texttt{uni2ts} \texttt{PackedStdScaler} contains an in-place
tensor assignment that breaks the autograd graph:
\texttt{loc[sample\_id == 0] = 0}.
We replace this with a differentiable \texttt{torch.where} operation,
enabling gradient-based fine-tuning through the scaler.
This fix is contributed back to the \texttt{uni2ts} community.

\textbf{Combined training objective.}
We fine-tune Moirai's parameters $\theta$ on a mixture of
70\% benign, 20\% hard-negative, and 10\% real-attack sequences
using the combined loss:
\begin{equation}
  \label{eq:combined_loss}
  \mathcal{L}(\theta) =
    \mathcal{L}_{\text{NLL}}(\theta) +
    \lambda\,\mathcal{L}_{\text{SupCon}}(\mathbf{z}, \mathbf{y};\, \tau),
\end{equation}
where $\lambda = 0.5$ and $\tau = 0.07$.
$\mathcal{L}_{\text{NLL}}$ is the standard negative log-likelihood loss
from Moirai's pre-training objective.
$\mathcal{L}_{\text{SupCon}}$~\cite{khosla2020supcon} is supervised
contrastive loss applied to embeddings $\mathbf{z}$ extracted from a
projection head $h_\phi$.

\textbf{Projection head.}
We insert a three-layer MLP projection head after Moirai's encoder:
$h_\phi: \R^{384} \to \R^{256} \to \text{ReLU} \to \R^{128}$.
Embeddings are L2-normalised before the contrastive loss, following
the recommendation of~\cite{khosla2020supcon}.

\textbf{Training details.}
We train with AdamW ($\eta = 10^{-4}$, weight decay $10^{-2}$),
batch size 32, for 10 epochs, with early stopping (patience 3) and
gradient clipping at 1.0.  We use the Moirai-Small (5M parameters)
checkpoint as the pre-trained backbone.
